% Research fragment prepared by /root/rectangular_bridge in run CUE_GCT-20261121.
% Requires amsmath, amssymb, amsthm; no nonstandard macros are used.
% Citations use bibliographic keys listed in sources.md.
% These are checked deductions in this run, not assertions of literature priority.

\section{Rectangular characters, explicit positive witnesses, and orbit ranks}

Throughout this section the ground field is $\mathbb C$, compact-group Haar
measure has mass one, and $D=nd$. We write $r=(d^n)$ and
$D_\alpha(b)=s_\alpha(1^b)=\dim S_\alpha(\mathbb C^b)$.
Duals are displayed when coordinate rings are involved. Passing from the
polynomial representation to the coordinate representation changes both sides
by the same duality and leaves every multiplicity below unchanged.

\subsection{R-01: the exact rectangular restriction transform}

\paragraph{Statement R-01.}
Let $a\geq\ell(\lambda)$ and let $X=(x_1,\ldots,x_a)$,
$Y=(y_1,\ldots,y_n)$. Denote the product alphabet by
$XY=(x_i y_j)_{1\leq i\leq a,1\leq j\leq n}$. Then
\begin{equation}\label{R01:branching}
 s_r(XY)=\sum_{\mu,\nu\vdash D}
 g(r,\mu,\nu)s_\mu(X)s_\nu(Y).
\end{equation}
In particular, by symmetry of the three Kronecker labels,
\begin{equation}\label{R01:haar}
 g(\lambda,r,r)=\int_{U(a)\times U(n)}
 s_r(A\otimes B)\overline{s_\lambda(A)}\det(B)^{-d}\,dA\,dB.
\end{equation}
Every coefficient needed on the rectangular-Kronecker side of determinant
GCT is therefore an explicit tensor-subgroup Fourier coefficient of the
\emph{same rectangular character} that occurs in the finite CUE moment.
The equality is not an identification of that coefficient with the scalar
moment.

More explicitly, enumerate the alphabet $XY$ as $z_1,\ldots,z_{an}$, with
$z_1,\ldots,z_n=x_1y_1,\ldots,x_1y_n$. Bump--Gamburd's Proposition~4
(source label \texttt{cpaverage}) gives
\begin{equation}\label{R01:CUE}
 s_r(XY)=\left(\prod_{j=1}^n z_j^d\right)
 \int_{U(d)}\prod_{j=1}^n\det(1+z_j^{-1}U^{-1})
 \prod_{j=n+1}^{an}\det(1+z_jU)\,dU.
\end{equation}
This formula is first read with nonzero $z_j$ and then as a polynomial identity.
Its specialization at identities proves the positive weighted sum
\begin{equation}\label{R01:dimension}
 D_r(an)=\sum_{\mu,\nu\vdash D}
 g(r,\mu,\nu)D_\mu(a)D_\nu(n)
\end{equation}
and the explicit upper bound
\begin{equation}\label{R01:bound}
 0\leq g(\lambda,r,r)\leq
 \left\lfloor\frac{D_r(an)}{D_\lambda(a)}\right\rfloor,
 \qquad
 D_r(an)=\prod_{i=1}^{n}\prod_{j=1}^{d}
 \frac{an+j-i}{n+d-i-j+1}.
\end{equation}
The numerator is a finite rectangular-character moment, not an unknown GCT
coefficient.

\begin{proof}
Apply Schur--Weyl duality to $(\mathbb C^a\otimes\mathbb C^n)^{\otimes D}$
in two ways. Decomposing first for $GL_{an}\times S_D$ gives the factors
$S_\eta(\mathbb C^a\otimes\mathbb C^n)\otimes[\eta]$.
Decomposing separately for $GL_a$ and $GL_n$ gives
$S_\mu(\mathbb C^a)\otimes S_\nu(\mathbb C^n)\otimes
([\mu]\otimes[\nu])$, with diagonal $S_D$ action. Comparing the multiplicity
space of $[r]$ proves \eqref{R01:branching}. Schur orthogonality and
$s_r(B)=\det(B)^d$ prove \eqref{R01:haar}. Equation \eqref{R01:CUE} is the
specialization $N=d,L=n,K=(a-1)n$ of the cited product identity.
Set every $x_i,y_j$ equal to one in \eqref{R01:branching}. Every summand is
nonnegative. Retain the term $(\mu,\nu)=(\lambda,r)$ and use
$D_r(n)=1$; this proves \eqref{R01:bound}. The displayed product is the
hook-content formula in Bump--Gamburd, equation \texttt{e\_hhl} in \S2.2.
\end{proof}

\paragraph{Finite moments and their limiting constants.}
Bump--Gamburd, Corollary~1, gives
\[
 \mathbb E_{U(N)}|\det(1-U)|^{2k}=D_{(N^k)}(2k).
\]
The limit after division by $N^{k^2}$ is
$c_k=\prod_{j=0}^{k-1}j!/(j+k)!$. Thus the finite moment is an integer
dimension, whereas the leading constant generally is not: $c_2=1/12$.
The distinction prevents a scalar asymptotic constant from being treated as
a representation multiplicity.

\subsection{R-02: swaps, tableaux, and positive projectors}

\paragraph{Statement R-02.}
Set $m=an$, $\delta_b=(b-1,b-2,\ldots,0)$, and
$\Delta_b(X)=\det(x_i^{b-j})$. For a tableau $T$ of shape $r$ with
semistandard entries in the ordered alphabet $\{1,\ldots,a\}\times
\{1,\ldots,n\}$, let $p(T)$ and $q(T)$ be the two marginal content vectors.
Write $N_r(p,q)$ for the number of these tableaux with marginal contents
$(p,q)$, and set it equal to zero when a content component is negative.
Then
\begin{align}\label{R02:tableaux}
 g(\lambda,r,r)
 &= [X^{\lambda+\delta_a}Y^{r+\delta_n}]
       \Delta_a(X)\Delta_n(Y)s_r(XY)\\
 &=\sum_{\sigma\in S_a,\tau\in S_n}
    \operatorname{sgn}(\sigma)\operatorname{sgn}(\tau)
    N_r(\lambda+\delta_a-\sigma\delta_a,
        r+\delta_n-\tau\delta_n).\nonumber
\end{align}
The exact unitary sum over swaps acts on the multiplicity by inserting
\begin{equation}\label{R02:swaps}
 s_{(d^n)}(z_1,\ldots,z_m)=
 \sum_{\substack{S\subseteq\{1,\ldots,m\}\\|S|=n}}
 \frac{\prod_{i\in S}z_i^d}
      {\prod_{i\in S,j\notin S}(1-z_j/z_i)}
\end{equation}
in the first line of \eqref{R02:tableaux}, and then setting $z_{ij}=x_i y_j$.
The sum in \eqref{R02:swaps} is combined before specialization or coefficient
extraction; its individual summands have poles and are not separately
nonnegative tableau counts.

There is, nevertheless, an exact positive formula. On the unitary space
\[
 T=S_\lambda(\mathbb C^n\otimes\mathbb C^n)
       \otimes\det(\mathbb C^n)^{-d}\otimes\det(\mathbb C^n)^{-d},
\]
let $P=\int_{U(n)\times U(n)}\rho_T(A,B)\,dA\,dB$. For any orthonormal
basis $(e_i)$ of $T$,
\begin{equation}\label{R02:projector}
 g(\lambda,r,r)=\operatorname{tr}P
   =\|P\|_{\mathrm{HS}}^2=\sum_{i,j}|\langle e_i,Pe_j\rangle|^2.
\end{equation}
In particular, this gives a rigorous sum of squares and a nonvanishing
criterion. It does not give a polynomial-size counting rule: the entries
of $P$ are themselves signed integrals, and the representation dimension
can be exponential in the natural parameters.

\begin{proof}
Multiply \eqref{R01:branching} by the two Vandermonde alternants.
For a partition $\mu$, the alternant $\Delta_a(X)s_\mu(X)$ has exactly one
monomial with strictly decreasing exponent vector equal to
$\lambda+\delta_a$ precisely when $\mu=\lambda$, and its coefficient then
is one. The analogous statement holds for $Y$. This proves the first line
of \eqref{R02:tableaux}. Expanding $s_r$ into semistandard tableaux and
the two alternants into their permutation terms proves the second line.
Formula \eqref{R02:swaps} is the rectangular Weyl-character Laplace expansion
used in Bump--Gamburd, Proposition~4 and \S4. For the last assertion,
Haar invariance gives $P^2=P$ and $P^*=P$; its range is exactly $T^{U(n)\times
U(n)}$. The tensor-subgroup branching identity identifies the dimension of
that range with $g(\lambda,r,r)$. Taking traces of $P^*P=P$ proves
\eqref{R02:projector}.
\end{proof}

\paragraph{What the arithmetic CFKRS recipe supplies.}
CFKRS's unitary formula is a theorem (Theorem~1.5.2 in the supplied source),
and its swap/contour conversion is Lemma~2.5.1. The corresponding
$L$-function mean-value assertion is Conjecture~1.5.1, derived by the
recipe in \S2.1 and the sum over $\Xi$ in \S2.2, equation
\texttt{conj:M}. Replacing the unitary factors in \eqref{R02:swaps} by the
arithmetic Euler product and shifted zeta factors does not define a
character restriction map. Neither RH nor GRH supplies the discarded
off-diagonal estimates in that recipe. The proven action on GCT is
\eqref{R02:tableaux}--\eqref{R02:swaps}; the arithmetic recipe remains a
separate conjectural mean-value statement.

\subsection{R-03: three explicit limits on positivity transfer}

\paragraph{Statement R-03.}
The following three proposed simplifications are false.
First, positivity of a rectangular CUE moment does not imply positivity
of every rectangular Kronecker coefficient. Second, the connected
left-right stabilizer multiplicity need not equal the full determinant
stabilizer multiplicity. Third, one cannot replace tensor-subgroup
branching by one multiplicity-free interlacing step.

\begin{proof}
For $r=(2,2)$ the $S_4$ representation $[r]$ is the two-dimensional
nontrivial summand in the permutation representation on the three
perfect matchings of four letters. Its character at the cycle types
$1^4,21^2,2^2,31,4$ is respectively $2,0,2,-1,0$; the class sizes are
$1,6,3,8,6$. The standard character $[31]$ is $3,1,-1,0,-1$.
Their character inner product gives
\[
 g((3,1),(2,2),(2,2))=(12-12)/24=0,
\]
although $D_{(2,2)}(4)=20>0$. The same elementary calculation yields
\[
 [22]\otimes[22]=[4]\oplus[22]\oplus[1111],\qquad
 \operatorname{Sym}^2[22]=[4]\oplus[22].
\]
Thus $g((1^4),(22),(22))=1$ but the corresponding symmetric Kronecker
coefficient is zero. These values are checked independently by enumerating
the 24 permutations in \texttt{check\_small\_cases.py}.

The tensor homomorphism $U(a)\times U(n)\longrightarrow U(an)$ has kernel
$\{(zI_a,z^{-1}I_n):|z|=1\}$. Its image has one-dimensional connected
center, whereas the block-diagonal Levi $U(a)\times U(n)$ has
two-dimensional connected center. For $a,n\geq2$ these are different
branching problems. Moreover, R-07 below constructs two independent
highest-weight invariants for $g((12,4),(4^4),(4^4))$, so its value is at
least two (in fact exactly two). A single $0/1$ interlacing rule cannot
be its multiplicity formula.
\end{proof}

The $0/1$ Gan--Gross--Prasad periods cited in Alp\"oge's steps~4 and~7
therefore require an additional, explicit comparison of branching
functors before they can control these multiplicities. Iterating
multiplicity-free restrictions can create multiple paths; no
contradiction with such restrictions is asserted. Likewise the squared
Gram formulas are not statements that every Weingarten summand is
positive: Collins--\'Sniady, Corollary~2.4 and its following example, give
$\operatorname{Wg}_N((12))=-1/(N(N^2-1))$ for $N\geq2$.

\subsection{R-04: plethysm requires the joint law of powers}

\paragraph{Statement R-04.}
For $N\geq\ell(\lambda)$ and $|\lambda|=nd$, the plethysm multiplicity is
\begin{equation}\label{R04:plethysm}
 a_\lambda(d[n])=\int_{U(N)}\overline{s_\lambda(U)}
 \sum_{\kappa\vdash d}\frac{1}{z_\kappa}
       \prod_{j\geq1}s_{(n)}(U^j)^{m_j(\kappa)}\,dU,
 \qquad z_\kappa=\prod_j j^{m_j(\kappa)}m_j(\kappa)!.
\end{equation}
Rains's Theorem~1.3 determines the eigenvalue law of each single $U^j$,
but replacing the different powers in \eqref{R04:plethysm} by independent
matrices with those marginal laws is invalid. A quantitative example is
\begin{equation}\label{R04:counter}
 \int_{U(2)}(\operatorname{tr}U)^4
        \overline{(\operatorname{tr}U^2)^2}\,dU=2,
\end{equation}
whereas the product of the two marginal expectations is zero.

\begin{proof}
For any finite-dimensional $W$, the cycle-index formula for the
symmetrizer $d!^{-1}\sum_{\sigma\in S_d}\sigma$ gives
$\chi_{\operatorname{Sym}^dW}(g)=\sum_{\kappa\vdash d}
z_\kappa^{-1}\prod_j\chi_W(g^j)^{m_j(\kappa)}$.
Set $W=\operatorname{Sym}^n\mathbb C^N$ and apply Schur orthogonality.
For \eqref{R04:counter}, in two variables
\[
 p_1^4=s_{(4)}+3s_{(3,1)}+2s_{(2,2)},\qquad
 p_2^2=s_{(4)}-s_{(3,1)}+2s_{(2,2)}.
\]
Orthogonality gives $1-3+4=2$. Multiplication of $U$ by a scalar phase
shows that both separate expectations vanish. The two polynomial
expansions are also verified by the exact computation certificate.
\end{proof}

In representation-ring language $U\mapsto U^j$ yields the Adams
operation on characters, not a homomorphism of nonabelian groups.
Already $\psi^2(V)=\operatorname{Sym}^2V-\bigwedge^2V$ has a negative
Schur coefficient. A positive alternative retaining the complete
coupling is
\begin{equation}\label{R04:wreath}
 a_\lambda(d[n])=\dim[\lambda]^{S_n\wr S_d}
 =\sum_{i,j}\left|\frac{1}{(n!)^d d!}
       \sum_{h\in S_n\wr S_d}\rho_\lambda(h)_{ij}\right|^2.
\end{equation}
Indeed, apply Schur--Weyl duality to $V^{\otimes nd}$ and the orthogonal
projector onto the wreath-product invariants, exactly as in R-02.
The formula is positive and exact; its factorial enumeration cost is
explicit, and it is not a solution of the efficient positive-formula
problem.

\subsection{R-05: Haar covariance computes the orbit-closure ranks}

\paragraph{Statement R-05.}
Let $V=\mathbb C^N$, $W=\operatorname{Sym}^nV$, $0\ne f\in W$, and
$X_f=\overline{GL(V)f}$. Put $E_d=\operatorname{Sym}^dW$ and
$v=f^{\otimes d}\in E_d$. Define
\begin{equation}\label{R05:twirl}
 M_f=\int_{U(N)}|\rho(k)v\rangle\langle\rho(k)v|\,dk.
\end{equation}
Under the Hermitian identification of $E_d$ with its conjugate dual,
$\ker M_f$ is the degree-$d$ ideal of $X_f$, and the image is dual to
$\mathbb C[X_f]_d$. Write the isotypic decomposition
\[
 E_d=\bigoplus_{\lambda\vdash nd} L_\lambda\otimes\mathbb C^{a_\lambda},
 \qquad L_\lambda=S_\lambda V,
 \qquad v_\lambda=\sum_{i=1}^{D_\lambda(N)}e_i\otimes m_i.
\]
Then the $\lambda$ block of $M_f$ is
\begin{equation}\label{R05:block}
 \frac{I_{L_\lambda}}{D_\lambda(N)}\otimes R_\lambda(f),
 \qquad R_\lambda(f)=\sum_i|m_i\rangle\langle m_i|.
\end{equation}
The multiplicity of $L_\lambda^*$ in $\mathbb C[X_f]_d$ is therefore
\begin{equation}\label{R05:rank}
 m_\lambda(f)=\operatorname{rank}R_\lambda(f).
\end{equation}
If $F_\lambda$ is the matrix with columns $m_i$, then, for every $s$,
\begin{equation}\label{R05:minors}
 e_s(R_\lambda(f))=
 \sum_{\substack{I\subseteq\{1,\ldots,a_\lambda\}\\|I|=s}}
 \ \sum_{\substack{J\subseteq\{1,\ldots,D_\lambda(N)\}\\|J|=s}}
       |\det(F_{\lambda,I,J})|^2.
\end{equation}
In particular $m_\lambda(f)\geq s$ if and only if this nonnegative
polynomial is positive. A multiplicity obstruction between two orbit
closures is thus an exact comparison of these Gram ranks, including
the boundary equations, rather than of their scalar traces.

\begin{proof}
For $u\in E_d$,
$\langle u,M_fu\rangle=\int|\langle u,\rho(k)v\rangle|^2\,dk$.
This is zero exactly when the continuous nonnegative integrand is
identically zero. The compact group $U(N)$ is Zariski dense in
$GL_N(\mathbb C)$: a regular function vanishing on it has all derivatives
along its real Lie algebra zero, and the complexification of that Lie
algebra is $\mathfrak{gl}_N(\mathbb C)$; analytic continuation on the
connected complex group gives zero. Consequently the same vanishing
holds on $GL(V)f$ and its closure. A degree-$d$ polynomial on $W$ is
exactly a linear functional on $E_d$ evaluated at $f^{\otimes d}$.
This identifies the kernel with the homogeneous ideal and the image
with the dual quotient.

Off-diagonal isotypic blocks of the twirl vanish by Schur's lemma.
Within $L_\lambda\otimes\mathbb C^{a_\lambda}$, averaging
$|e_i\rangle\langle e_j|$ gives
$\delta_{ij}I_{L_\lambda}/D_\lambda(N)$, proving
\eqref{R05:block}. Thus the orbit span is
$L_\lambda\otimes\operatorname{span}\{m_i\}$, which proves
\eqref{R05:rank}. Finally $R_\lambda=F_\lambda F_\lambda^*$.
The coefficient $e_s$ is the sum of its principal $s$-minors;
Cauchy--Binet gives \eqref{R05:minors}. Since a positive semidefinite
matrix has nonnegative eigenvalues, $e_s>0$ exactly when at least $s$
of those eigenvalues are positive.
\end{proof}

For $f=\det_n$ or the padded permanent, \eqref{R05:twirl} is a Haar
integral over $U(n^2)$. Its entries have degree $nd$ in the matrix
entries and degree $nd$ in their conjugates. Thus the exact
Collins--\'Sniady formula evaluates them. The full, dimension-truncated
Weingarten inverse must be used when tensor degree exceeds group
dimension. The remaining task is to control the ranks of the resulting
structured matrices; knowing that the matrices are positive
semidefinite does not determine which minors vanish.

\subsection{R-06: a Gaussian and integral form of the rank certificate}

\paragraph{Statement R-06.}
Use the tensor realization $E_d\subseteq V^{\otimes D}$ and set
$v=f^{\otimes d}$. For an $N\times N$ matrix of indeterminates $Z$, write
\begin{equation}\label{R06:coefficients}
 Z^{\otimes D}v=\sum_{|\alpha|=D}Z^\alpha c_\alpha,
 \qquad C_f=(c_\alpha)_\alpha.
\end{equation}
There are $N^D$ tensor coordinates and
$\binom{N^2+D-1}{D}$ columns. If $G$ has independent standard complex
Gaussian entries, then
\begin{equation}\label{R06:gaussian}
 \mathbb E|G^{\otimes D}v\rangle\langle G^{\otimes D}v|
   =\sum_{|\alpha|=D}\alpha!\,c_\alpha c_\alpha^*.
\end{equation}
Its image is the orbit span from R-05. It has the same kernel and
the same isotypic ranks as the Haar covariance.

Let $f^\lambda=\dim[\lambda]$ and let
\[
 e_\lambda=\frac{f^\lambda}{D!}
       \sum_{\sigma\in S_D}\chi^\lambda(\sigma)\sigma
\]
act by permutation of tensor factors. Then
\begin{equation}\label{R06:integer}
 A_\lambda(f)=\sum_{\sigma\in S_D}\chi^\lambda(\sigma)\sigma C_f,
 \qquad
 \operatorname{rank}A_\lambda(f)=D_\lambda(N)m_\lambda(f).
\end{equation}
When $f$ is an integral symmetric tensor, $A_\lambda(f)$ has integer
entries. If its tensor entries have absolute value at most $H$, then
every entry of $A_\lambda(f)$ has absolute value at most
$D!f^\lambda N^D H^d$.
The same rank is obtained by replacing the coefficient columns by all
evaluations at the finite grid $Z\in\{0,1,\ldots,D\}^{N^2}$ and then
applying $e_\lambda$. This grid and the coefficient matrix are
explicit but not polynomial-size constructions.
An ordinary polynomial with integer monomial coefficients may require
scaling by $n!$ under the symmetric-tensor embedding to meet the
integrality hypothesis; this nonzero scalar does not change any orbit
span or multiplicity in characteristic zero.

\begin{proof}
For a standard complex Gaussian $z$, polar integration gives
$\mathbb E[z^p\bar z^q]=0$ if $p\ne q$ and $p!$ if $p=q$.
Independence proves \eqref{R06:gaussian}. Its image is
$\operatorname{span}\{c_\alpha\}$. A linear functional kills all
$c_\alpha$ exactly when its evaluation on $Z^{\otimes D}v$ is the zero
polynomial. Since invertible matrices are Zariski dense, that is
equivalent to killing the $GL_N$ orbit span. This proves the image and
kernel assertions without analytic approximation.

The central idempotent $e_\lambda$ projects precisely onto the
$L_\lambda$ isotypic summand by Schur--Weyl duality. Consequently
$e_\lambda C_f$ has rank $D_\lambda(N)m_\lambda(f)$, and multiplication
by the nonzero scalar $D!/f^\lambda$ gives \eqref{R06:integer}.
Symmetric-group characters and permutation matrices have integral
entries. Each tensor coordinate of $Z^{\otimes D}v$ is a sum of $N^D$
terms whose coefficients have magnitude at most $H^d$; hence
$|(c_\alpha)_i|\leq N^D H^d$. The character bound
$|\chi^\lambda(\sigma)|\leq f^\lambda$ and the $D!$ permutations give
the stated estimate. Finally a polynomial of individual degree at
most $D$ vanishing on a Cartesian product of $D+1$ distinct scalars is
zero, by induction on the number of variables. Applying this fact to
every linear functional proves equality of the grid and coefficient
spans.
\end{proof}

This replaces signed Haar expansions by an explicitly positive
coefficient Gram matrix. Cauchy--Binet gives positive sums of squared
integer minors after clearing the indicated denominators. It also
exhibits the exact cost: the tensor model has $N^D$ rows, the
symmetrizer has $D!$ terms, and the displayed grid has
$(D+1)^{N^2}$ points. A claimed efficient positivity algorithm must
compress these data or exploit their ranks; positivity alone does
not achieve that compression.

\subsection{R-07: explicit positive witnesses that survive the boundary}

\paragraph{Statement R-07.}
Let $n\geq2$, $d\geq0$, and $0\leq q\leq d$. For
$\lambda=(nd-q,q)$, define
\[
 p_{2,\ldots,n}(q)=[t^q]\prod_{j=2}^n(1-t^j)^{-1}.
\]
There are explicitly constructed, linearly independent highest-weight
vectors, indexed by partitions of $q$ into parts $2,\ldots,n$, in the
degree-$d$ coordinate ring of the determinant orbit closure. In fact
\begin{equation}\label{R07:positive}
 m_\lambda(\det_n)=a_\lambda(d[n])=p_{2,\ldots,n}(q),
 \qquad
 g(\lambda,(d^n),(d^n))\geq p_{2,\ldots,n}(q).
\end{equation}
The last inequality is equality by Manivel, Corollary~2(1) in \S3.3; it is a specialization of his
stable rectangular Kronecker Theorem~1, restated as
Ikenmeyer--Panova, Theorem~2.1. The construction below gives concrete witnesses for this
known numerical positivity and, simultaneously, proves their survival
on the determinant orbit closure.

More generally, with no restriction $q\leq d$,
\begin{equation}\label{R07:alltwo}
 m_\lambda(\det_n)=a_\lambda(d[n])
 \quad\text{for every }\lambda\vdash nd\text{ with }\ell(\lambda)\leq2.
\end{equation}
Therefore no multiplicity obstruction of length at most two can
separate a padded-permanent orbit closure from the determinant orbit
closure in the same ambient space and degree.

\begin{proof}
Write a binary degree-$n$ form as
$F(s,t)=\sum_{j=0}^n\binom nj b_j s^{n-j}t^j$.
In the coordinate representation $\operatorname{Sym}^n\mathbb C^2$,
the upper-triangular raising operator is
$E b_j=j b_{j-1}$, with $E b_0=0$, and the weight of $b_j$ is
$(n-j,j)$. For $2\leq j\leq n$ put
\begin{equation}\label{R07:central}
 Q_j=b_0^{j-1}\sum_{i=0}^j\binom ji b_i
                  \left(-\frac{b_1}{b_0}\right)^{j-i}.
\end{equation}
This is a polynomial: for $i=0$ the factor $b_i=b_0$ cancels the only
possible negative exponent; all $i\geq1$ terms already have
nonnegative powers of $b_0$. The first two examples are
\[
 Q_2=b_0b_2-b_1^2,\qquad
 Q_3=b_0^2b_3-3b_0b_1b_2+2b_1^3.
\]
Under $\exp(uE)$, $b_i$ goes to
$\sum_{h=0}^i\binom ih u^{i-h}b_h$ and $b_1/b_0$ goes to
$b_1/b_0+u$. The binomial theorem in \eqref{R07:central} shows that
the two translations cancel. Thus $E Q_j=0$. Every term has degree
$j$ in the $b_i$ and weight $((n-1)j,j)$.

For a tuple of nonnegative integers $(a_2,\ldots,a_n)$ satisfying
$\sum_{j=2}^n ja_j=q$, define
\begin{equation}\label{R07:witness}
 P_{\mathbf a}=b_0^{d-q}\prod_{j=2}^n Q_j^{a_j}.
\end{equation}
This is a highest-weight polynomial of degree $d$ and weight
$(nd-q,q)$. Evaluate it on determinant pencils
$F(s,t)=\det(sI_n+tY)$, with $Y$ diagonal. Every monic binary form
is obtained in this way over $\mathbb C$ by factoring the corresponding
one-variable polynomial. In particular, one may set $b_0=1,b_1=0$
and choose $b_2,\ldots,b_n$ arbitrarily. On this family $Q_j=b_j$,
so the restrictions of the $P_{\mathbf a}$ are distinct monomials
in independent variables. They are linearly independent.

For completeness, the ambient highest-weight dimension in this range
is obtained by a short weight count. In
$\operatorname{Sym}^d(\operatorname{Sym}^n\mathbb C^2)$ the weight
$(nd-q,q)$ has dimension equal to the number of partitions of $q$
into at most $d$ parts, each at most $n$. When $q\leq d$ this is
$[t^q]\prod_{j=1}^n(1-t^j)^{-1}$. A $GL_2$ irreducible of highest
weight $(nd-r,r)$ has one vector at each weight $(nd-q,q)$ with
$r\leq q\leq nd-r$. Subtracting consecutive weight dimensions
therefore gives $p_{2,\ldots,n}(q)$. This proves that the displayed
independent witnesses form a basis and establishes the first equality
in \eqref{R07:positive} once boundary survival is checked.

Now let $V$ have dimension $n^2$ and view a binary form as an element
of $\operatorname{Sym}^nV^*$ supported on the first two coordinates.
It factors as $\prod_{i=1}^n\ell_i$ and hence equals
$\det(\operatorname{diag}(\ell_1,\ldots,\ell_n))$.
This is a linear, possibly singular, substitution into $\det_n$.
Since $GL(V)$ is dense in $\operatorname{End}(V)$ and substitution
is polynomial, every such form lies in the determinant orbit closure.
A vector of highest weight $\lambda=(\lambda_1,\lambda_2,0,\ldots,0)$
in $\operatorname{Sym}^d(\operatorname{Sym}^nV)$ uses only the first
two basis vectors: each other weight coordinate is a sum of
nonnegative exponents and equals zero. The additional raising
operators annihilate it, so its highest-weight condition is exactly
the binary one. A nonzero polynomial in these binary coordinates
cannot vanish on every binary form. Thus the homogeneous ideal of
the determinant orbit closure has no nonzero highest-weight vector
of length at most two. Taking the quotient proves
\eqref{R07:alltwo}, including all degrees and all two-row weights.

The determinant-coordinate multiplicity is at most its orbit
multiplicity, which is the symmetric rectangular Kronecker coefficient,
and that is at most $g(\lambda,r,r)$; this is the cited stabilizer
statement in B\"urgisser--Ikenmeyer--Panova, \S1(b).
This proves the inequality in \eqref{R07:positive}. Finally any
padded-permanent coordinate ring is a quotient of the same ambient
$\operatorname{Sym}(\operatorname{Sym}^nV)$, so its multiplicity is
at most $a_\lambda(d[n])=m_\lambda(\det_n)$ for two-row $\lambda$.
\end{proof}

For $n=d=q=4$, the two witnesses are $Q_4$ and $Q_2^2$, with
\[
 Q_4=b_0^3b_4-4b_0^2b_1b_3+6b_0b_1^2b_2-3b_1^4.
\]
They specialize to $b_4$ and $b_2^2$, proving directly that the
determinant coordinate multiplicity at $(12,4)$ is two. This is
stronger information than nonnegativity of a Haar integral.

\subsection{R-08: the exact Molien numerator needed for orbit closures}

\paragraph{Statement R-08.}
Let $S=\operatorname{Sym}(W)$ be a polynomial coordinate algebra,
$A=S/I$ a graded $GL_N$-equivariant quotient, and $K=U(N)$. Then
\begin{equation}\label{R08:molien}
 \sum_{d\geq0}m_\lambda(A_d)t^d
   =\int_K\overline{\chi_\lambda(k)}H_A(t,k)\,dk.
\end{equation}
For the ambient algebra,
$H_S(t,k)=\det(1-t\rho_W(k))^{-1}$.
For a finite graded equivariant free resolution
$\cdots\to S\otimes T_1\to S\otimes T_0\to A\to0$,
\begin{equation}\label{R08:numerator}
 H_A(t,k)=
 \frac{\sum_i(-1)^i\chi_{T_i}(t,k)}
      {\det(1-t\rho_W(k))}.
\end{equation}
Thus the Molien denominator and ratios identities compute the desired
multiplicities only after the orbit ideal's equivariant numerator,
or equivalently the ranks in R-05--R-06, is supplied.

\begin{proof}
Coefficientwise Schur orthogonality proves \eqref{R08:molien}.
The eigenvalue expansion of the symmetric algebra gives the
denominator formula. Alternating graded characters are additive
on exact complexes, yielding \eqref{R08:numerator}. For a purely
formal interpretation no convergence assumption is needed; for
analytic evaluation take $|t|<1$ in the unitary polynomial case.
\end{proof}

\subsection{R-09: an exact ratios kernel for squared GCT multiplicities}

\paragraph{Statement R-09.}
Let $n,d,N\geq1$, $r=(d^n)$, and define
\[
 S_N(n,d)=\sum_{\substack{\lambda\vdash nd\\\ell(\lambda)\leq N}}
                    g(\lambda,r,r)^2.
\]
Set $S_0(n,d)=0$.
For $0<t<1$ and $A,B,C,D\in U(n)$ put
\begin{align*}
 \mathcal R_N(t;A,B,C,D)=\int_{U(N)}
 &\det\bigl(I-t(A\otimes B)\otimes U\bigr)^{-1}\\[-2mm]
 &\quad\cdot
 \det\bigl(I-t(C\otimes D)\otimes U^{-1}\bigr)^{-1}\,dU.
\end{align*}
Then the explicitly specified fourfold Fourier projection is
\begin{equation}\label{R09:squares}
 \int_{U(n)^4}\mathcal R_N(t;A,B,C,D)
   (\det A\det B\det C\det D)^{-d}\,dA\,dB\,dC\,dD
       =t^{2nd}S_N(n,d).
\end{equation}
For $N\geq n^2$, the inner ratio is evaluated exactly by
Bump--Gamburd, Theorem~3, at $K=L=0,Q=R=n^2$:
\begin{equation}\label{R09:evaluated}
 \mathcal R_N(t;A,B,C,D)
       =\det\bigl(I_{n^4}-t^2 A\otimes B\otimes C\otimes D\bigr)^{-1}.
\end{equation}
Consequently an explicitly evaluated Molien--Weyl coefficient is the
sum of squares of precisely the rectangular Kronecker coefficients
needed in GCT. The truncated sums satisfy
\begin{equation}\label{R09:monotone}
 S_N(n,d)-S_{N-1}(n,d)
    =\sum_{\substack{\lambda\vdash nd\\\ell(\lambda)=N}}g(\lambda,r,r)^2
    \geq0,
 \qquad
 g(\lambda,r,r)\leq\lfloor\sqrt{S_N(n,d)}\rfloor
 \quad(\ell(\lambda)\leq N).
\end{equation}
The stabilized value has a second nonnegative moment expression:
\begin{equation}\label{R09:finitefourth}
 S_{n^2}(n,d)=\frac1{(nd)!}\sum_{\sigma\in S_{nd}}\chi^r(\sigma)^4.
\end{equation}

\begin{proof}
The Cauchy identity and absolute convergence for $t<1$ expand the two
inverse determinants into
$\sum_\lambda t^{|\lambda|}s_\lambda(A\otimes B)s_\lambda(U)$
and
$\sum_\mu t^{|\mu|}s_\mu(C\otimes D)s_\mu(U^{-1})$.
Schur orthogonality on $U(N)$ retains exactly $\lambda=\mu$ of
length at most $N$. The two applications of the tensor-subgroup
projection in R-02 then retain $|\lambda|=nd$ and give
$g(\lambda,r,r)$ from each pair of $U(n)$ factors. This proves
\eqref{R09:squares}. When $N\geq n^2$, the length cutoff is
inactive because both product alphabets have $n^2$ entries. A second
Cauchy identity gives \eqref{R09:evaluated}, equivalently the stated
specialization of the ratios theorem. The definition immediately
proves \eqref{R09:monotone}. Finally the character of
$[r]\otimes[r]$ is $(\chi^r)^2$. Its squared character norm is
the sum of squares of its irreducible multiplicities, proving
\eqref{R09:finitefourth}; symmetric-group characters are real.
The tensor-subgroup branching formula also gives
$g(\lambda,r,r)=0$ for $\ell(\lambda)>n^2$, because
$S_\lambda(\mathbb C^n\otimes\mathbb C^n)=0$ in that range.
Thus the cutoff in the left side of \eqref{R09:finitefourth}
omits no nonzero multiplicity.
\end{proof}

The positive quantities in \eqref{R09:squares} are obtained after
projection. The pointwise ratio and individual swap summands need
not be nonnegative. An approximation argument would need a
certified error below one half to recover an integer coefficient,
with bounds uniform in the growing tensor degree. Neither the
closed ratio formula nor RH supplies those estimates or compresses
the remaining fourfold Fourier projection.

\subsection{R-10: the complete height-two rectangle and its boundary gate}

\paragraph{Statement R-10 (known formula, specialized and compared).}
Let $d\geq0$ and $\lambda\vdash2d$. For partitions of length at
most four, append zero parts to obtain $(\lambda_1,\ldots,\lambda_4)$.
Then
\begin{equation}\label{R10:rectangle}
g(\lambda,(d,d),(d,d))=
\begin{cases}
1,&\ell(\lambda)\leq4\text{ and all four parts have the same parity},\\
0,&\text{otherwise}.
\end{cases}
\end{equation}
More precisely, the multiplicity of $[\lambda]$ in
$\operatorname{Sym}^2[d,d]$ is one exactly for partitions with all
parts even and length at most four. Its multiplicity in
$\bigwedge^2[d,d]$ is one exactly for partitions with exactly four
odd parts. These formulas are the constant-term specialization of
Blasiak--Mulmuley--Sohoni, GCT~IV, Corollary~18.4
(source label \texttt{c kronecker generating function sym ext});
they are cited known results, not new numerical formulas.

In the determinant problem with $n=2$, $N=4$, the boundary side is
also completely determined:
\begin{equation}\label{R10:closure}
m_\lambda(\det_2)=a_\lambda(d[2])=
\begin{cases}
1,&\ell(\lambda)\leq4\text{ and every part of }\lambda\text{ is even},\\
0,&\text{otherwise}.
\end{cases}
\end{equation}
Thus all four-odd-part summands of the ordinary rectangular
coefficient are absent from the determinant coordinate ring.
For $d=2$, this includes the $\lambda=(1,1,1,1)$ example in R-03.
Simultaneous conjugation of the two inner partitions also extends
\eqref{R10:rectangle}, as an ordinary Kronecker formula, to rectangles
of width two.

\begin{proof}
Write $c_i=\lambda_i-\lambda_{i+1}$ for $1\leq i\leq3$ and
$c_4=\lambda_4$, so $c_i$ is the number of columns of height $i$.
GCT~IV, Corollary~18.4, gives the generating function
$\sum_{\alpha\vdash2d,\ell(\alpha)\leq2}
g_{\varepsilon,\alpha,\lambda}x^{\alpha_1-\alpha_2}$.
Three of its four cases have a factor $x$ or $x^2$ multiplying
polynomials with nonnegative exponents, and hence contribute no
constant term. The remaining case is the product
$\llbracket c_1\rrbracket\llbracket c_2\rrbracket
\llbracket c_3\rrbracket$, where
$\llbracket b\rrbracket=x^b+x^{b-2}+\cdots$ ends at $1$ or $x$.
Its constant coefficient is one exactly when $c_1,c_2,c_3$ are all
even; the case condition then requires
$\varepsilon=(-1)^{c_3+c_4}=(-1)^{\lambda_4}$.
Thus the symmetric square retains the all-even partitions and the
exterior square retains the four-odd-part partitions. Their sum is
\eqref{R10:rectangle}. Length greater than four gives zero by
tensor-subgroup branching. Finally
$[\rho']=[\rho]\otimes\operatorname{sgn}$, so conjugating both
inner partitions leaves their tensor product unchanged.

The determinant $x_{11}x_{22}-x_{12}x_{21}$ is a nondegenerate
quadratic form in four variables. Over $\mathbb C$ all nondegenerate
quadratic forms of a fixed dimension are congruent. They form a
nonempty Zariski-open subset of the space of quadrics, so
$\operatorname{Det}_2=\operatorname{Sym}^2(\mathbb C^4)^*$.
Consequently its coordinate ring is
$\operatorname{Sym}(\operatorname{Sym}^2\mathbb C^4)$.

Here is a direct highest-weight proof of its familiar multiplicity
formula. In $\operatorname{Sym}(\operatorname{Sym}^2\mathbb C^s)$
write the generators as the entries of a symmetric matrix $Q$,
with $q_{ij}$ of weight $e_i+e_j$. The upper-triangular action on
these generators is $Q\mapsto g^tQg$.
Let $\Delta_i$ be the leading principal $i\times i$ minor and
$\Delta_0=1$. Each $\Delta_i$ is fixed by the upper unitriangular
group and has weight $2(e_1+\cdots+e_i)$.
On the dense open set where all $\Delta_i$ are nonzero, elementary
upper-triangular congruences reduce $Q$ to the unique diagonal
matrix with entries $\Delta_i/\Delta_{i-1}$.
Therefore a nonzero unitriangular-invariant polynomial has nonzero
restriction to diagonal matrices. If it is also a torus eigenvector
of weight $\lambda$, this restriction is a scalar multiple of
$\prod_i q_{ii}^{\lambda_i/2}$. It exists only when each
$\lambda_i$ is even, and its value determines the polynomial on a
dense open set. Hence each highest-weight space has dimension at
most one and vanishes unless $\lambda=2\mu$.
Conversely, for every partition $\mu$ of length at most $s$,
\[
 \prod_{i=1}^s\Delta_i^{\mu_i-\mu_{i+1}},\qquad\mu_{s+1}=0,
\]
is a nonzero highest-weight polynomial of weight $2\mu$ and
degree $|\mu|$. This proves
$\operatorname{Sym}^d(\operatorname{Sym}^2\mathbb C^s)
=\bigoplus_{\mu\vdash d,\ell(\mu)\leq s}S_{2\mu}\mathbb C^s$
and, at $s=4$, \eqref{R10:closure}.
\end{proof}

The stabilizer interpretation agrees with these formulas.
The full linear stabilizer of the nondegenerate quadratic form is
$O_4$, whereas the connected tensor image of $SL_2\times SL_2$ is
$SO_4$: its kernel is $\{(I,I),(-I,-I)\}$, so its dimension is six,
the dimension of $SO_4$. The extra orthogonal component removes
the odd sector. Thus the connected-stabilizer upper bound, the
symmetric rectangular coefficient, and the closure multiplicity
are completely distinguished in this case.

\subsection{R-11: small-degree padding excludes all four-row obstructions}

\paragraph{Statement R-11.}
Work over $\mathbb C$. Let $1\leq m\leq3$, $n\geq m$, $N=n^2$,
$V=\mathbb C^N$, $f\in\operatorname{Sym}^mV^*$, and $\ell\in V^*$.
Set $h=\ell^{n-m}f$. For every coordinate-ring degree $d\geq0$
and every $\lambda\vdash nd$ with $\ell(\lambda)\leq4$,
\begin{equation}\label{R11:fourrow}
 m_\lambda(h)\leq m_\lambda(\det_n).
\end{equation}
The assertion includes padded permanents of degree at most three.
It is a deduction from classical linear determinantal
representations, not a claim that arbitrary $f$ itself has an
$n\times n$ determinantal representation.

\begin{proof}
First consider forms in at most four variables. Linear forms are
$1\times1$ determinants. Every quadric in at most four variables
is congruent over $\mathbb C$ to $y_1^2+\cdots+y_r^2$, $r\leq4$.
With unused $y_i$ set to zero it is the determinant
\begin{equation}\label{R11:quadric}
 \det\begin{pmatrix}
 y_1+\mathrm i y_2&y_3+\mathrm i y_4\\
 -y_3+\mathrm i y_4&y_1-\mathrm i y_2
 \end{pmatrix}=
 y_1^2+y_2^2+y_3^2+y_4^2.
\end{equation}
Beauville (2000), Corollary~6.4, states over an algebraically
closed field that every smooth cubic surface has a $3\times3$
linear determinantal equation. Thus over $\mathbb C$ every smooth
quaternary cubic is a $3\times3$ linear determinant, after
absorbing its nonzero scalar factor in one matrix row.
Smooth cubics are Zariski dense in the affine space of quaternary
cubics: the discriminant-complement is open and nonempty, as the
Fermat cubic shows. Therefore every quaternary cubic lies in the
closure of the linear-determinant image.

If $q=\det M$ has degree $m$ and $u$ is linear, then
$u^{n-m}q=\det\operatorname{diag}(M,uI_{n-m})$.
The entries of this block matrix give a linear substitution into
$\det_n$. Such substitutions belong to $\operatorname{Det}_n$,
because invertible linear maps are dense in all linear maps and
substitution is polynomial. For fixed $u$, the preimage of
$\operatorname{Det}_n$ under the linear map $q\mapsto u^{n-m}q$
is closed. The density just proved extends the assertion from
smooth cubics to every cubic. We have shown that
\begin{equation}\label{R11:projectedcontainment}
 u^{n-m}q\in\operatorname{Det}_n
 \quad\text{for every }q\text{ of degree }m\leq3
 \text{ and }u\text{ supported on at most four variables}.
\end{equation}

Put $k=\min(4,N)$ and let $\pi$ set the last $N-k$ variables
equal to zero. Let $P$ be a highest-weight vector of weight
$\lambda$ in the degree-$d$ ideal of $\operatorname{Det}_n$.
Since $\lambda_j=0$ for $j>k$ and all coordinate generators have
nonnegative weights, every monomial of $P$ uses only coefficients
of forms supported on the first $k$ variables. Hence
$P(F)=P(\pi F)$ for every $F\in\operatorname{Sym}^nV^*$.
For $g\in GL(V)$ the projection is multiplicative on forms, so
\[
 \pi(g\cdot h)=(\pi(g\cdot\ell))^{n-m}\pi(g\cdot f)
 \in\operatorname{Det}_n
\]
by \eqref{R11:projectedcontainment}. It follows that
$P(g\cdot h)=0$ for every $g$, and hence $P$ belongs to the
ideal of $\overline{GL(V)\cdot h}$.
Thus the highest-weight space of the determinant ideal is
contained in that of the padded-form ideal for every such
$\lambda$. Finite-dimensional polynomial $GL(V)$-modules are
semisimple in characteristic zero. Subtracting the ideal
multiplicity from the common ambient multiplicity gives
\eqref{R11:fourrow}. If $N<4$, partitions of length greater than
$N$ already have multiplicity zero, so the same statement holds.
\end{proof}

\subsection{R-12: exact permanent-pencil dominance and four-row deficits}

\paragraph{Statement R-12 (computer-assisted finite-degree deduction).}
For $n=4,5$, the polynomial map
\[
 \Phi_n:(\operatorname{Mat}_n\mathbb C)^4
 \longrightarrow\operatorname{Sym}^n(\mathbb C^4)^*,
 \qquad(A_1,A_2,A_3,A_4)\longmapsto
 \operatorname{per}\!\left(\sum_{k=1}^4x_kA_k\right)
\]
is dominant. Consequently, for every $d\geq0$ and
$\lambda\vdash nd$ with $\ell(\lambda)\leq4$,
\begin{equation}\label{R12:perambient}
 m_\lambda(\operatorname{per}_n)=a_\lambda(d[n])
 \qquad(n=4,5).
\end{equation}
Combining the quartic case with G-07 and writing $D=320112$ gives
the exact comparison
\begin{equation}\label{R12:deficit}
 m_\lambda(\operatorname{per}_4)-m_\lambda(\det_4)
 =a_{\lambda-(D,D,D,D)}((d-D)[4]),
\end{equation}
with the same zero conventions as in G-07.
In particular, for every integer $s\geq0$,
\begin{equation}\label{R12:gapone}
 d=D+s,\qquad \lambda=(D+4s,D,D,D)
 \quad\Longrightarrow\quad
 m_\lambda(\operatorname{per}_4)-m_\lambda(\det_4)=1.
\end{equation}
These are multiplicity obstructions; no occurrence-vanishing
claim is implied. For $n=5$, G-06 and \eqref{R12:perambient}
give a four-row multiplicity obstruction in some degree, without
specifying that degree or highest weight.

\paragraph{Explicit rank certificate for $n=4$.}
Each entry of the following array is the coefficient quadruple of
a linear form in the ordered variables $(x_1,x_2,x_3,x_4)$:
\begin{equation}\label{R12:pencil4}
\begin{pmatrix}
(2,-1,0,0)&(2,-2,0,-2)&(-2,-1,-1,0)&(0,-1,-1,2)\\
(0,1,0,0)&(2,-2,-1,1)&(-2,1,-2,0)&(-1,-2,-1,-1)\\
(0,-2,0,2)&(1,-2,-1,-1)&(2,-1,-1,0)&(2,2,0,0)\\
(1,1,2,1)&(1,0,2,-2)&(-1,2,-1,1)&(2,1,1,1)
\end{pmatrix}.
\end{equation}
Call the resulting matrix of linear forms $L_4(x)$. Index the
parameters $a_{ijk}$ by $0\leq i,j\leq3$, $0\leq k\leq3$,
where $a_{ijk}$ multiplies $x_{k+1}$ in entry $(i+1,j+1)$,
and assign column number $c=4(4i+j)+k$.
Order the 35 exponent quadruples $\alpha$ of total degree four in
decreasing lexicographic order. Define the integer Jacobian by
\begin{equation}\label{R12:derivative}
 J_{\alpha,,4(4i+j)+k}
 =[x^\alpha]\;x_{k+1}
 \operatorname{per}\bigl(L_4(x)_{\widehat{i+1},\widehat{j+1}}\bigr).
\end{equation}
Take all 35 rows and the ordered columns
\[
 C_4=(0,1,\ldots,30,32,33,34,35).
\]
The resulting square integer matrix has determinant
\begin{equation}\label{R12:minor4}
 \det J_{C_4}
 =45459683955923995466096193465479594741071872
 \equiv77\pmod{101}.
\end{equation}
Thus \eqref{R12:pencil4}--\eqref{R12:minor4} specify the entire
nonvanishing certificate in the text, independently of a file format.

\paragraph{Explicit rank certificate for $n=5$.}
Use the same coefficient-quadruple convention with the array
\begin{equation}\label{R12:pencil5}
\begin{pmatrix}
(2,2,-1,1)&(2,-2,-2,1)&(-2,-1,1,1)&(1,-1,0,0)&(1,2,0,2)\\
(1,0,0,-1)&(1,2,0,2)&(1,2,-1,2)&(0,2,-1,1)&(-1,2,1,0)\\
(-2,-2,1,0)&(1,-1,-1,2)&(0,1,0,-1)&(2,2,-1,1)&(2,-1,1,-2)\\
(2,0,0,1)&(-1,1,2,2)&(0,0,1,-2)&(1,-1,1,-1)&(-1,0,1,-2)\\
(0,-2,1,-1)&(1,0,-1,1)&(-1,-1,1,-1)&(-1,-1,2,1)&(0,-2,-2,2)
\end{pmatrix}.
\end{equation}
There are 56 degree-five coefficient rows, again in decreasing
lexicographic order, and parameter column $4(5i+j)+k$.
Use the same derivative rule, with $4\times4$ complementary
permanents, and the columns
\[
 C_5=(0,1,\ldots,37,39,40,\ldots,56).
\]
The selected $56\times56$ minor has determinant $30$ modulo 101.

\begin{proof}
Differentiating the permanent with respect to one matrix entry
gives the complementary permanent; differentiating that entry
with respect to $a_{ijk}$ contributes $x_{k+1}$. This proves the
Jacobian rule. Expanding the complementary permanents and
performing exact modular elimination gives the displayed
nonzero residues. For full reproducibility,
\texttt{check\_permanent\_projection.py} records the full pencils,
coefficient matrices, selected columns, and integer minors.
The characteristic-zero determinants are also calculated by
fraction-free Bareiss elimination. As a second check on the
derivative expansion, every one of the 64 quartic columns was
recomputed as the exact finite difference
$\Phi_4(a+e_{ijk})-\Phi_4(a)$; this is exactly the derivative
because the map is affine in each individual parameter.
Twelve quintic columns were checked in the same way.

A nonzero minor modulo 101 is a nonzero integer minor, so the
complex differential is surjective at each displayed point.
Fix the other parameters and retain the selected 35 or 56
parameters. The resulting square polynomial map has invertible
complex differential there. By the inverse function theorem its
image contains a nonempty analytic open subset. A polynomial
vanishing on that image vanishes on an analytic open subset and
is identically zero. This proves Zariski density, hence dominance.

Every permanent pencil is a linear endomorphism substitution
into $\operatorname{per}_n$, and therefore belongs to its
$GL_{n^2}$ orbit closure. This closed variety consequently
contains the full linear space of degree-$n$ forms supported on
the first four variables. A highest-weight polynomial of weight
$\lambda$ with $\ell(\lambda)\leq4$ uses only those coefficient
variables, by the nonnegative-weight support argument of R-07.
If it vanishes on the permanent orbit closure, it vanishes on
that entire coefficient space, so it is the zero polynomial.
Thus the ideal has no highest-weight vector in those weights,
which proves \eqref{R12:perambient}.

Equation \eqref{R12:deficit} now follows by subtracting G-07.
For \eqref{R12:gapone}, the shifted partition is $(4s)$, whose
multiplicity in $\operatorname{Sym}^s(\operatorname{Sym}^4V)$ is
one: its weight space is spanned by $(e_1^4)^s$, and this vector
is highest weight. This also covers $s=0$.
For $n=5$, the four-variable determinant image has dimension at
most $52<56$ by G-06, so its ideal contains a nonzero homogeneous
highest-weight vector. G-05 excludes highest weights of length
at most three. Hence it removes at least one four-row ambient
multiplicity, whereas \eqref{R12:perambient} retains them all.
\end{proof}

These computations are exact finite certificates. The independent
agent recheck uses Ryser's formula at integer evaluation points
and an invertible interpolation matrix, so it does not reuse
the complementary-permanent coefficient-expansion algorithm.
The checks establish the stated degrees four and five; they
do not extrapolate dominance to growing $n$ or supply an
asymptotic permanent-versus-determinant lower bound.

\subsection{R-13: the exact range of quaternary permanent dominance}

\paragraph{Statement R-13 (finite certificates and a dimension obstruction).}
For the polynomial map $\Phi_n$ in R-12,
\begin{equation}\label{R13:range}
 \Phi_n\text{ is dominant over }\mathbb C
 \quad\Longleftrightarrow\quad1\leq n\leq16.
\end{equation}
Consequently, for $2\leq n\leq16$ and every
$\lambda\vdash nd$ with $\ell(\lambda)\leq4$,
\begin{equation}\label{R13:multiplicities}
 m_\lambda(\operatorname{per}_n)=a_\lambda(d[n]).
\end{equation}
For every $4\leq n\leq16$, the unpadded permanent and determinant
of size $n$ have a multiplicity obstruction with exactly four rows
in some coordinate-ring degree. R-12 gives an explicit infinite
family when $n=4$; no degree or partition is asserted here for
$5\leq n\leq16$. This is a certificate-based finite-range result,
with no claim of literature priority or asymptotic separation.

\paragraph{Complete specification of the finite certificates.}
All arithmetic for this paragraph is in $\mathbb F_{101}$, with
integer coefficient representatives in $\{0,\ldots,100\}$.
For each $n=1,\ldots,16$ set
\[
 s_0=20261121+1009n,\qquad
 s_{c+1}=(1664525s_c+1013904223)\bmod2^{32},\qquad
 a_c=\left\lfloor\frac{s_{c+1}}{2^{16}}\right\rfloor\bmod101.
\]
There are $4n^2$ coefficients, with $0\leq c<4n^2$, and the
pencil is
\[
 L_n(x)_{i+1,j+1}=\sum_{k=0}^3a_{4(ni+j)+k}x_{k+1}
 \quad(0\leq i,j<n).
\]
Put $b_n=\binom{n+3}{3}$ and order all integer points
$\alpha\in\mathbb N^4$ with $|\alpha|=n$ in decreasing
lexicographic order. Define the $b_n\times4n^2$ matrix
\begin{equation}\label{R13:evaluationmatrix}
 E_n\bigl(\alpha,4(ni+j)+k\bigr)
 =\alpha_{k+1}\operatorname{per}
   \bigl(L_n(\alpha)_{\widehat{i+1},\widehat{j+1}}\bigr)
 \pmod{101}.
\end{equation}
Choose columns greedily from left to right, retaining a column
exactly when it increases the span of previously retained columns.
Equivalently, scan columns left to right in Gaussian elimination,
choose the first available nonzero pivot row, and clear the rows
below it. This produces $b_n$ columns $C_n$ in every case, and the
square minors, with rows and columns in the stated order, have
the following nonzero residues:
\[
\begin{array}{c|r|r@{\qquad}c|r|r}
n&b_n&\det(E_n)_{C_n}&n&b_n&\det(E_n)_{C_n}\\\hline
1&4&1&9&220&67\\
2&10&11&10&286&67\\
3&20&64&11&364&42\\
4&35&49&12&455&87\\
5&56&85&13&560&25\\
6&84&28&14&680&69\\
7&120&25&15&816&73\\
8&165&59&16&969&8
\end{array}
\]
These uniform range certificates use the displayed generator;
R-12 retains its separate, smaller-entry quartic and quintic pencils.
The directory \texttt{permanent\_range} contains all generated
integer coefficients, evaluation points, chosen columns, and full
matrices, with unambiguous row-major byte encoding and checksums.
Thus the formula above and the code reproduce the certificates
without probabilistic choices or an external data source.

\begin{proof}
For the positive direction, \eqref{R13:evaluationmatrix} is the
derivative of the coefficient map followed by evaluation at the
specified points, by the complementary-permanent derivative
rule in R-12. Its full row rank $b_n$ forces the coefficient-map
differential, whose target dimension is $b_n$, to be surjective.
The nonzero minor modulo 101 is the reduction of a nonzero
integer minor. Hence the characteristic-zero differential has
the same required full rank, and the inverse-function argument
in R-12 proves dominance.

Here are the exact algorithms used to certify every entry and
minor in the table. For a numeric $n\times n$ matrix $M$, let
$F(\varnothing)=1$ and
\[
 F(S)=\sum_{j\in S}M_{|S|,j}F(S\setminus\{j\}).
\]
Induction on $|S|$ shows that $F(S)$ is the permanent formed from
the first $|S|$ rows and columns $S$. Thus $F([n])$ is the
permanent. Reverse differentiation of this recurrence computes
every entry derivative, hence every complementary permanent,
in $O(n2^n)$ field operations per evaluation point. This is
implemented in \texttt{check\_permanent\_range.cpp}.

All $3{,}348{,}864$ entries across the sixteen matrices were then
recomputed from Ryser's different identity
\[
 \operatorname{per}(M)=
 \sum_{S\subseteq[n]}(-1)^{n-|S|}
 \prod_{i=1}^n\left(\sum_{j\in S}M_{ij}\right).
\]
To prove this identity, expand the products as maps
$f:[n]\to[n]$. The coefficient of a fixed map is
$\sum_{S\supseteq\operatorname{im}f}(-1)^{n-|S|}$, which is zero
unless $f$ is onto and is one if it is onto. Such onto maps are
exactly permutations. Differentiating this identity and using
Gray-code row-sum updates gives the independent algorithm in
\texttt{check\_permanent\_range\_ryser.cpp}. Every entry agrees
with the first computation. Finally
\texttt{check\_permanent\_range\_minors.py} recomputes each
selected determinant by a second modular elimination algorithm;
all sixteen nonzero residues agree. These are finite exact
checks, rather than floating-point rank estimates.

For the converse, the torus
\[
 H_n=\{(u_1,\ldots,u_n,v_1,\ldots,v_n)\in(\mathbb C^*)^{2n}:
                  \prod_i u_i\prod_j v_j=1\}
\]
preserves $\Phi_n$ under the entrywise action
$L_{ij}\mapsto u_i v_jL_{ij}$. It has dimension $2n-1$.
On the nonempty open set where every entry $L_{ij}$ is a nonzero
linear form, a stabilizing torus element satisfies
$u_i v_j=1$ for all $i,j$. It is therefore of the form
$u_i=t$, $v_j=t^{-1}$, a one-dimensional stabilizer.
Thus each generic fiber contains an orbit of dimension $2n-2$,
and
\begin{equation}\label{R13:dimension}
 \dim\overline{\operatorname{im}\Phi_n}
       \leq4n^2-2n+2.
\end{equation}
At $n=17$, this is $1124<1140=b_{17}$. For $n\geq18$,
\[
 6\bigl(b_n-(4n^2-2n+2)\bigr)
 =n^3-18n^2+23n-6>0.
\]
Consequently no $n\geq17$ map is dominant. This proves
\eqref{R13:range}.

For $n\geq2$, the ambient variable space has dimension $n^2\geq4$,
so dominance places the entire space of degree-$n$ forms in the
first four variables inside the permanent orbit closure. The
highest-weight support argument of R-12 proves
\eqref{R13:multiplicities}.
For $n\geq4$, G-06 gives a nonzero determinant equation supported
on four variables, and G-05 excludes weights with at most three
rows. Thus some exactly four-row highest weight has determinant
multiplicity below its ambient value. Equation
\eqref{R13:multiplicities} retains that ambient value on the
permanent side when $n\leq16$, giving the asserted obstruction.
\end{proof}

The dimension ceiling concerns full ambient survival in four
variables. It does not determine the comparison between the two
proper four-variable images for $n\geq17$.

\subsection{R-14: a dimension criterion that survives genuine padding}

\paragraph{Statement R-14.}
Let $\mathcal P_{m,4}$ be the Zariski closure of the image of
$\Phi_m$ and put $c_m=\dim\mathcal P_{m,4}$.
For integers $n>m\geq2$, embed the usual padded permanent
$h_{m,n}=z^{n-m}\operatorname{per}_m$ in
$\operatorname{Sym}^n(\mathbb C^{n^2})^*$, with $z$ independent
of the $m^2$ permanent variables. The projection of its orbit
closure to the first four variables is exactly
\begin{equation}\label{R14:paddingvariety}
 \mathcal Z_{m,n}=
 \{\ell^{n-m}q:\ell\in(\mathbb C^4)^*,\ q\in\mathcal P_{m,4}\},
 \qquad \dim\mathcal Z_{m,n}=c_m+3.
\end{equation}
If
\begin{equation}\label{R14:criterion}
 c_m+3>2n^2+2,
\end{equation}
then, for every sufficiently large coordinate-ring degree $d$,
there is a partition $\lambda\vdash nd$ of length exactly four
such that
\[
 m_\lambda(h_{m,n})>m_\lambda(\det_n).
\]
No effective threshold for $d$ is asserted.

For $4\leq m\leq16$, R-13 gives $c_m=\binom{m+3}{3}$, so the
sufficient condition becomes
\begin{equation}\label{R14:explicitrange}
 m<n\leq\left\lfloor\sqrt{\frac12\binom{m+3}{3}}\right\rfloor.
\end{equation}
The nonempty padded ranges are
\[
\begin{array}{c|ccccccccc}
m&8&9&10&11&12&13&14&15&16\\\hline
n&9&10&11&12\text{--}13&13\text{--}15&14\text{--}16&
15\text{--}18&16\text{--}20&17\text{--}22
\end{array}.
\]
In particular, some four-row multiplicity obstruction separates
the unpadded $\operatorname{per}_n$ and $\det_n$ for every
$4\leq n\leq22$; for $17\leq n\leq22$ this follows by
block-diagonal containment from the $m=16$ padded cases.

\begin{proof}
First, $\mathcal P_{m,4}$ is an irreducible affine cone, as the
closure of the image of an irreducible affine space. It contains
an irreducible degree-$m$ polynomial: use the cyclic matrix with
diagonal entries $x_1$, entries $(i,i+1)$ equal to $x_2$ for
$i<m$, entry $(m,1)$ equal to $x_3$, and every other entry zero.
Its only perfect matchings are the diagonal and the full cycle,
so its permanent is $x_1^m+x_2^{m-1}x_3$.
Viewed in $\mathbb C[x_2,x_3][x_1]$, this polynomial is
Eisenstein at the prime $x_3$ and hence irreducible.
The reducible degree-$m$ forms constitute a closed subset:
projectivize each of the finitely many multiplication maps of
degrees $a$ and $m-a$, $1\leq a<m$, and take their closed images.
It follows that a general $q\in\mathcal P_{m,4}$ is irreducible.

Projectivizing the multiplication map
$(\ell,q)\mapsto\ell^{n-m}q$ gives a morphism
\[
 \mathbb P^3\times\mathbb P(\mathcal P_{m,4})
 \longrightarrow\mathbb P(\operatorname{Sym}^n\mathbb C^4)^*.
\]
Its image is closed because its domain is projective; its affine
cone is exactly the set in \eqref{R14:paddingvariety}, including
zero. For a generic irreducible $q$ of degree $m\geq2$, the
linear factor $\ell$ is uniquely determined by the product,
up to scalar, and then so is $q$. Thus the projective map is
generically injective. Its image dimension is
$3+(c_m-1)$ and the affine-cone dimension is $c_m+3$.

Every four-variable restriction of an invertible substitution
in $h_{m,n}$ has the form in \eqref{R14:paddingvariety}, so
closedness gives one inclusion of the projected orbit closure.
Conversely, for each permanent pencil $q=\Phi_m(A)$, any chosen
$\ell$ and its $m^2$ matrix-entry linear forms can be prescribed
simultaneously by an endomorphism substitution, since the
padding variable and the permanent variables are independent.
There is enough room because $m^2+1\leq n^2$.
Such substitutions lie in the orbit closure. Taking limits in
$q$ then shows that the whole closed set $\mathcal Z_{m,n}$ is
contained in the orbit closure, supported on its first four
variables. This proves equality of the projected image.

Let $\mathcal D_{n,4}$ be the determinant four-variable image
closure. By G-06, $\dim\mathcal D_{n,4}\leq2n^2+2$.
The coordinate rings of both cones are standard graded. The
Hilbert function of a nonzero standard graded affine cone of
dimension $c$ agrees for large $d$ with a polynomial of degree
$c-1$ and positive leading coefficient. Thus
\eqref{R14:criterion} implies
\[
 \dim\mathbb C[\mathcal Z_{m,n}]_d
 >\dim\mathbb C[\mathcal D_{n,4}]_d
\]
for every sufficiently large $d$.
Decompose these finite-dimensional $GL_4$ representations into
irreducibles. Their dimensions are the sums of multiplicities
times the positive integers $\dim S_\lambda\mathbb C^4$.
Hence some irreducible multiplicity on the first side is larger.

These $GL_4$ multiplicities agree with the corresponding
$GL_{n^2}$ multiplicities of length at most four in the original
orbit-closure coordinate rings. Indeed, highest-weight vectors
with those weights use only the first-four coefficient variables,
and the intersection of the full ideal with that polynomial
subring is the ideal of its projected image. The additional
raising operators involving later variables vanish on such
vectors. This proves the claimed multiplicity inequality.
G-05 states that the determinant multiplicity is already the
ambient maximum for weights of length at most three, so the
strict inequality must occur at length exactly four.

Substitution of $c_m=\binom{m+3}{3}$ into
\eqref{R14:criterion} gives $2n^2\leq\binom{m+3}{3}$, which is
\eqref{R14:explicitrange} and the displayed finite table.
Finally
\[
 \operatorname{per}\operatorname{diag}(L_m,zI_{n-m})
 =z^{n-m}\operatorname{per}(L_m)
\]
shows that the padded-permanent orbit closure is contained in
the size-$n$ permanent orbit closure. The resulting coordinate
quotient is surjective and therefore cannot increase
multiplicities. The padded gaps at $m=16$ thus imply the stated
unpadded gaps for $17\leq n\leq22$.
\end{proof}

If the expected formula
$c_m=\min\{\binom{m+3}{3},4m^2-2m+2\}$ were proved beyond the
certified range, \eqref{R14:criterion} would turn it into further
padded obstructions up to approximately $n=\sqrt2\,m$.
This expectation is a separate generic-fiber problem; neither
the finite certificates nor the dimension upper bound proves it.
The criterion above is unconditional whenever its stated
dimension lower bound has been established.

\subsection{What the attached GCT IV, VII, and VIII positivity claims supply}

The attached GCT~IV is the later Blasiak--Mulmuley--Sohoni manuscript
corresponding to their 2015 memoir, with filename
\texttt{Apr13KroneckerGCT4.tex}. Its \S12.1 theorem
(\texttt{t nonstandard schur-weyl duality}) proves nonstandard
Schur--Weyl duality for general $\dim V,\dim W$; the later
explicit-basis limitation should not obscure this proved
structural result. Its Theorem~15.21(vii)--(viii)
(\texttt{t main canonical basis}) proves a highest-weight crystal
count for $g(\alpha,\beta,\gamma)$ when
$\ell(\alpha),\ell(\beta)\leq2$, including the symmetric/exterior
split. Theorem~18.2 (\texttt{t consequence for Kronecker}) makes
the count explicit, and Corollary~18.4 gives the generating
function used in R-10. This supplies a positive rule for rectangles
of height two, or width two after simultaneous conjugation.
The construction for the full tensor power remains Conjecture~19.1
(\texttt{cj canonical basis X\^{}r}); \S19.2 describes the failed
flatness of a naive extension beyond the two-row case.

Even in the proved case, counting highest-weight elements does not
mean that every matrix coefficient is nonnegative. GCT~IV,
Example~17.8 (\texttt{ex not positive}), displays the divided-power
coefficient $q^{10}-q^4-q^{-4}+q^{-10}$. This concerns the specified
basis and operator; it is a concrete warning against replacing a
precise canonical-basis positivity statement by unrestricted
coefficient positivity.

GCT~VII (Mulmuley, 2008a, revised extended abstract) proves in
Theorem~1.2 (\texttt{tnonstd}) existence of the nonstandard
quantization, complete reducibility of its polynomial
representations, and its Peter--Weyl statement. The nonstandard
duality and reciprocity needed to recover plethysm are explicitly
Conjectures~1.3--1.4 (\texttt{cduality}, \texttt{creciprocity});
the refined form is Conjecture~6.1 (\texttt{crecirefined}).
The proposed formula
$a_{\lambda,\mu}^{\pi}=\sum_\alpha
m_\lambda^\alpha n_\pi^\alpha$ is conditional on those structures.
The introduction states that correctness of the basis algorithms
and their required properties has not been proved there.

GCT~VIII (Mulmuley, 2008b, revised extended abstract) makes the
additional gates explicit: local crystal existence is Conjecture~2.2
(\texttt{clocalcrystalb}); minimum-degree globalization is
Conjectures~2.5, 3.1--3.2; cellular decomposition is Conjectures~2.6,
3.3 and the specialization-compatible quasi-cell condition is
Conjecture~3.7. Kronecker positivity is Conjectures~2.7 and~3.4,
with specified signs and powers of $q-q^{-1}$ allowed. The general
plethysm setting uses the weaker saturation and nonstandard
positivity Conjectures~2.8--2.9 and~3.5--3.6. Finally, \S6
(\texttt{scomplexity}) also requires efficient combinatorial
labels and crystal operations to obtain a $\#P$ formula; these
complexity properties are themselves proposed rather than proved.

Its discussion of the Riemann hypothesis in the introduction and
\S2.4.2 is a suggested extension of finite-field purity and
perverse-sheaf machinery to the nonflat, noncommutative objects at
issue. It is not a theorem deriving these conjectures from the
classical RH or Dirichlet GRH. Accepting the stipulated analytic
theorems therefore leaves specific algebraic gates: construct the
required realization, prove its cellular and specialization
properties, and control the sizes of its labels and operations.
The supplied papers prove the height-two case and propose this
general route; they supply no unconditional constructive positive
formula for arbitrary rectangles and no theorem controlling the
orbit-closure ranks in R-05--R-06.

\paragraph{Algebraic-geometric dictionary for R-01--R-08.}
R-01 and R-02 are identities of Schur operations and restriction maps
on representation $K_0$, evaluated over $\mathbb C$ for Haar integration.
The tensor homomorphism is the integral group-scheme morphism
$GL_a\times GL_n\to GL_{an}$, and its kernel contains the central
$\mathbb G_m$ described in R-03. R-04 is the cycle-index expression
for symmetric powers in a $\lambda$-ring; the Adams operations preserve
the representation ring, not its cone of actual representations.
For R-05--R-06 the intrinsic algebraic object is the graded orbit-map
image, or the coefficient map in \eqref{R06:coefficients}. Its rank
strata are defined by determinantal (Fitting) ideals. The Gaussian
and Hermitian metrics are analytic devices for the same rank; they
are not asserted to be motives. The Schur idempotent is defined over
$\mathbb Z[1/D!]$, and the integral matrix in \eqref{R06:integer}
avoids treating characteristic-zero semisimplicity as an integral
direct-sum theorem. R-07 is a statement about restrictions of locally
free Schur constructions to the binary-form subscheme and the
scheme-theoretic image of the linear-substitution morphism. Its
formulas can be cleared of their binomial denominators over
$\mathbb Z$, while the basis and orbit-closure claims here are in
characteristic zero. R-08 is additivity in equivariant $K_0$ applied
to a resolution; the numerator remembers the quotient's equations.
Exact SGA/EGA/FGA and cote~133 locators for this dictionary are supplied
in the separate Grothendieck part of the paper.
R-09 identifies a four-factor invariant of
$\operatorname{Sym}(V_1\otimes V_2\otimes V_3\otimes V_4)$ with the
composition of two tensor-product restriction functors; its degree
$nd$ determinant-isotypic coefficient is the fourfold Kronecker
multiplicity. The finite CUE cutoff is the rank restriction on a
Schur functor, not an arithmetic purity hypothesis.
